\section{The remaining route to complete recognition}\label{sec:recognition}

The proved algorithm completely decides a bounded simplification question.
To use it as a complete unknot recognizer, one must connect that question
to the topology of every relevant exterior. This section states a sufficient
condition precisely, so that the missing research target is not hidden
inside the phrase ``small barrier.''

\subsection{A conditional closure theorem}

Fix the maintained diagram-to-exterior construction and its source binding.
Suppose an $n$-crossing diagram produces, in polynomial time with a verified
source binding, a marked triangulation with at most $(n+2)^{O(1)}$
tetrahedra and polynomial marking bit length. Let $\mathcal C$
be a class containing these exteriors and closed under the permitted
simplifications. A \emph{sound terminal predicate} $P$ is a computable
predicate on members of $\mathcal C$ such that $P(T)$ implies that $T$ is a
solid torus. Its acceptance must be bound to the actual exterior; a
recognizer for an unrelated supplied triangulation is insufficient.

\begin{theorem}[A sufficient global structural target]\label{thm:conditional}
Assume the following two properties for the chosen exterior class and move
model.
\begin{enumerate}
\item The predicate $P$ is decidable in
$2^{O((\log(t+B+2))^2)}$ time and is sound for solid tori.
\item There is an absolute constant $C$ such that every solid torus
$T\in\mathcal C$ with $P(T)$ false admits a strict descent with at most
$C(\log(t+2))^2$ total upward moves.
\end{enumerate}
Then the bounded-descent algorithm and $P$ yield a deterministic
$n^{O(\log(n+2))}$ unknot-recognition algorithm for those input diagrams.
\end{theorem}
\begin{proof}
Let $t_0$ be the initial source size and choose the fixed integer budget
$U=\lceil C(\log(t_0+2))^2\rceil$. At each current triangulation first evaluate $P$.
If it accepts, accept the input. Otherwise run complete bounded-$U$ descent.
Adopt any independently verified positive endpoint and repeat. If the
bounded search exhausts, reject.

Every successful iteration strictly decreases the tetrahedron count, so
there are fewer than $t_0$ such iterations. All replacements preserve the
underlying manifold. If the input is knotted, its exterior is not a solid
torus, so soundness prevents acceptance by $P$. If the input is an unknot,
every current source is a solid torus in $\mathcal C$. Unless $P$ accepts,
the second hypothesis supplies a descent within the chosen budget, since
the current size is at most $t_0$. Complete exhaustion cannot then cause
rejection. Finite termination therefore accepts exactly the unknot inputs.

By \cref{cor:preprocess}, the cumulative marking bit bound increases by at
most $t_0U+O(1)$, remaining polynomial in $n$. Each descent call costs
$2^{O(U)}\poly(t_0+B)=n^{O(\log(n+2))}$, and the terminal calls have the
same form. Multiplying by a polynomial number of iterations preserves
that bound. The maintained exterior binding transfers the conclusion to
the original diagram.
\end{proof}

This theorem is a conditional reduction, not a proof of its second
hypothesis. It is useful because it permits \emph{any} successful bounded
descent at each iteration: the structural hypothesis applies to every
resulting solid torus in the closed class. A claim concerning only the
existence of one favorable long path would not justify the same greedy
choice of intermediate endpoints.

The predicate need not recognize all solid tori. It must cover the
solid-torus states that cannot be further simplified within the proposed
budget. A terminal family may consequently be much larger than a handful
of smallest triangulations. In this move model both the vertices and the
boundary triangulation are preserved. These invariants constrain the
reachable terminal states and must be part of any valid global theorem.

\subsection{Three distinctions that cannot be suppressed}

\paragraph{Total upward cost and maximum height.}
A trace can alternate many upward and downward moves while remaining close
to its initial tetrahedron count. A bound on the largest intermediate
triangulation therefore does not bound the total number of upward events.
The localization and counting theorems charge total upward cost. Converting
a height-barrier theorem into this stronger resource bound requires an
additional compression argument, potentially using causal components,
redundant loops, or a monotone potential. It is not obtained by changing
the notation for the parameter.

\paragraph{Descent and disc discovery.}
The component argument uses additive tetrahedron loss. A disc found after
several independent modifications can depend on the combined endpoint
geometry and cochain. Its existence need not be witnessed after projecting
just one positive component. Any analogous theorem for normal-disc search
needs a separate localization invariant for that topological endpoint
property. The implementation's shared-cover engine can still accelerate
that search, but its complete bounded region is then a supplied hypothesis.

\paragraph{Bounded irreducibility and knot type.}
A complete failure to descend within budget is an exact statement about
the allowed moves. It identifies neither a canonical triangulation nor a
nontrivial knot by itself. The terminal and structural hypotheses in
\cref{thm:conditional} are what turn bounded failure into a topological
decision. Resource-limited failure is weaker still and remains inconclusive.

\section{Further research questions and proposed milestones}\label{sec:questions}

The questions below separate structural mathematics from implementation
work. Each has a concrete target whose success would change the next
algorithmic decision.

\subsection{Structural questions with direct asymptotic consequences}

\paragraph{1. Find a terminal class admitting a small total-upward theorem.}
For the exact finite exterior class produced by ProveIt, identify a sound
terminal predicate $P$ and prove, or refute by a family, the second
hypothesis of \cref{thm:conditional}. Boundary triangulation and vertex data
must be retained in the statement. A useful intermediate result would
establish $U=O(\log^2 t)$ for a natural subclass such as a controlled family
of layered or hierarchy-derived solid tori. The milestone is a uniform
theorem over that subclass, with its constants and encoding costs explicit.

\paragraph{2. Relate height barriers to total upward cost.}
Let $h$ bound the excess tetrahedra along a simplification path. What extra
condition on that path, its dependency graph, or its hierarchy controls
the minimum total upward cost by $\poly(h,\log t)$? A negative family with
small $h$ and large minimum total upward cost would also be valuable: it
would rule out a tempting but invalid route to the required bound. Recording
both parameters in future searches is a necessary first experiment.

\paragraph{3. Determine the optimal localization function.}
Define $f(U)$ as the least universal radius guaranteeing a connected
first-descent footprint whenever descent with at most $U$ upward events
exists in the stated category. This report proves $f(U)\le3U+3$ and
$f(1)=6$. Is the linear coefficient three asymptotically sharp, or do
geometric constraints force extra birth arcs and improve the incidence
estimate? Any proposed lower-bound family must rule out \emph{all}
alternative low-cost simplifications, not only exhibit a large chosen
trace. The six-cell proof illustrates the required strength.

\paragraph{4. Generalize the additive objective carefully.}
Which other integer potentials admit a componentwise additive loss and a
bounded-per-event incidence census? A potential involving tetrahedra,
specific normal coordinates, or a hierarchy measure could make more
terminal states accessible without leaving the causal framework. The
target is an abstract theorem with explicit local input/output sizes,
backward persistence, and a positive-component argument. Nonadditive
normal-disc existence should first be tested for counterexamples to naive
projection.

\paragraph{5. Extend the move set without losing the proof contract.}
Boundary moves, vertex-changing moves, or $2$--$0$ reductions may overcome
the preserved invariants of the present model. Each addition changes the
legality and input/output census. One must establish whether every enabling
dependency still consumes a born cell, or whether a richer graph is needed.
The strict-simplicial suspension example gives a concrete diagnostic:
global missing-simplex tests create dependencies invisible to tetrahedron
births. A useful result would quantify localization for a larger,
precisely specified move system rather than appeal to Pachner connectivity
without resource bounds.

\subsection{Algorithmic questions suggested by the new theorems}

\paragraph{6. Complete the deterministic multiunit implementation.}
Implement a reproducible perfect-hash-family generator, stream authenticated
positive items through the supplied-coloring DP, reconstruct the selected
certificates, and call the disjoint-trace composer. An end-to-end negative
answer must certify that every required coloring and item family was
exhausted. The first milestone is agreement with unrestricted small-instance
search under a \emph{joint} upward budget. Enumerating components separately
and greedily repeating descents is not a substitute for this comparison.

\paragraph{7. Reduce index memory while preserving exact cover semantics.}
The full inverted index removes repeated search but stores all connected
covers. Can a compact decision diagram, canonical-parent search oracle,
or incremental connected-completion structure represent
$\{R:F\subseteq R\}$ exactly with substantially lower practical memory?
An admissibility oracle must distinguish actual footprints from their
possible connecting covers. It must also support resource-capped failure
without misreporting complete exhaustion. Compare streaming restart,
explicit index, and compact index on both positive and fully exhausted
instances, including setup time and peak memory.

\paragraph{8. Prove a stronger safe quotient of continuation states.}
The present deduplication saves endpoint queries but retains continuation
frames. A geometric isomorphism alone forgets original-source supports and
sleep history. Determine the smallest additional data making two frames
future-equivalent under the remaining upward budget and cover family.
A proof of a right congruence would justify further pruning. Until then,
endpoint equality remains evidence about outputs, not permission to merge
all histories.

\paragraph{9. Improve local geometry updates and marking storage.}
Global validation and rebuilding are polynomial but dominate small traces.
Maintain affected edge stars, dual incidences, and cochain transport using
stable local frames and persistent data structures. The acceptance contract
must remain independent: a faster producer should still emit certificates
checked without calling its own move or transport routines. Test every
accepted local frame, particularly downward belt relabelings, as part of
the update invariant rather than relying only on the producer's canonical
examples.

\paragraph{10. Tighten the trace-count constants and use budget structure.}
The $11$-symbol encoding deliberately permits many impossible commitments.
Can local compatibility, first-descent balance, or birth-graph sparsity give
a smaller exponential base without costly per-state reasoning? A second
parameter such as pathwidth of the consumed region may support a different
dynamic program. The milestone is a proved count and a measured crossover,
not a branching-factor estimate based only on successful fixtures.

\subsection{Evidence, verification, and integration questions}

\paragraph{11. Build a heterogeneous barrier census.}
Extend the current constructed solid-torus controls with genuine diagram
exteriors, independently recognized hard solid tori, local-minimum
triangulations, and adversarial cases from published Pachner-graph
exploration. Record total upward cost, maximum height, minimum actual
footprint, component structure, index size, peak memory, and complete or
capped status. Retain unsuccessful cases and certificate replay. Such a
census can guide conjectures about \cref{thm:conditional}; it cannot prove
the required asymptotic bound on its own.

\paragraph{12. Certify negative bounded search and formalize the core.}
Positive traces already have an independent geometric checker. Can complete
bounded failure be exported as a compact exploration DAG whose checker
verifies initial state, all outgoing alternatives, resource budgets,
commutation justifications, and terminal rejection? Separately, formalize
the small core consisting of backward persistence, component projection,
footprint connectivity, and the incidence theorem. A formalization should
start from explicit generalized face pairings and local permutations;
ordinary strict simplicial complexes have different legality rules. This
would give the repository a precise trust boundary for future optimizations.

\subsection{A practical order of work}

The immediate integration milestone is the additive cover search with
independent replay and the default radius $3U+3$. The next implementation
milestone is complete multiunit color coding and a broader measured census.
In parallel, the highest-value mathematical task is the terminal-class
and total-upward structural theorem. Better constants and data structures
make the current bounded primitive more useful, while that structural
theorem is what could close the project's specific quasi-polynomial
recognition argument.

\section{Deliverables and integration contract}\label{sec:deliverables}

The archive contains the complete article sources and compiled PDF, the
new producer and verifier modules, causal certificate utilities, packing
prototype, fixtures, targeted tests, exact raw benchmark records, positive
certificates, independent audit results, and a self-contained reproduction
snapshot. The snapshot combines the pinned native tree with only the
previously absent incoming dependencies and the additive files from this
report. It is not an instruction to overwrite the current repository with
the whole snapshot.

The integration script defaults to a dry run. It checks recorded native
prerequisite hashes and refuses conflicting destination files; applying
the overlay only creates absent files or accepts byte-identical existing
ones. A manifest identifies native baseline files, incoming dependencies,
and this report's additions separately. Reproduction scripts retain new
outputs separately from the recorded evidence. Checksums make the delivered
bytes auditable, while mathematical correctness rests on the proofs and
geometric verification rather than on those checksums.

The resulting contribution is a complete fixed-parameter simplification
theorem, its multiunit strengthening, and a measured implementation of
shared local search with independently checked witnesses. The project-wide
recognition consequence is stated with its remaining assumptions in
\cref{thm:conditional}, so it can be pursued as a precise research problem.
